% ============================================================
%  CAPÍTULO 16 – Vías Futuras
% ============================================================
\chapter{Vías Futuras}

FitLabs, en su versión actual, constituye un MVP sólido y funcional. Sin embargo, el potencial de crecimiento es significativo. A continuación se presentan las líneas de evolución más relevantes para una versión futura del producto:

\section{Mejoras técnicas}

\begin{itemize}
  \item \textbf{Implementación de un gestor de estado}: Migrar la gestión del estado de la aplicación a Riverpod o BLoC para obtener mayor predictibilidad, testabilidad y rendimiento en pantallas con alta frecuencia de actualización.
  \item \textbf{Pruebas unitarias y de widgets}: Incorporar una suite de tests automatizados con \texttt{flutter\_test} para los módulos críticos (autenticación, creación de rutinas).
  \item \textbf{Modo offline}: Implementar un mecanismo de caché local (SQLite o Hive) para permitir el uso básico de la aplicación sin conexión a Internet.
  \item \textbf{Notificaciones push}: Integrar Firebase Cloud Messaging (FCM) para notificar al entrenador sobre nuevos mensajes o próximas sesiones.
\end{itemize}

\section{Nuevas funcionalidades}

\begin{itemize}
  \item \textbf{Módulo de nutrición}: Seguimiento del plan nutricional del cliente, con registro de macros y calorías diarias.
  \item \textbf{Integración con wearables}: Sincronización automática con Google Fit, Apple Health o dispositivos Garmin/Fitbit para importar datos de actividad del cliente.
  \item \textbf{Videollamadas integradas}: Módulo de sesiones de entrenamiento online con videoconferencia, aprovechando APIs como Daily.co o Agora.
  \item \textbf{Sistema de pagos}: Integración con Stripe o PayPal para la gestión de suscripciones y pagos de sesiones directamente desde la app.
  \item \textbf{Panel web para entrenadores}: Versión web de la interfaz del entrenador (aprovechando que Flutter compila a Web) para gestión desde escritorio.
  \item \textbf{Recomendación automática de rutinas}: Sistema de sugerencia de rutinas basado en el histórico del cliente y sus objetivos, utilizando un modelo de Machine Learning ligero o la API de OpenAI.
\end{itemize}

\section{Escalabilidad y comercialización}

\begin{itemize}
  \item Migrar de Supabase Free a Supabase Pro (o auto-alojado) para eliminar las limitaciones del plan gratuito.
  \item Publicar la aplicación en Google Play Store y Apple App Store.
  \item Implementar un modelo de negocio \textit{freemium}: plan gratuito limitado a 5 clientes y plan profesional con clientes ilimitados y funcionalidades avanzadas.
\end{itemize}

\section{Roadmap propuesto a 12 meses}

Para transformar FitLabs en un producto comercialmente viable, se plantea una hoja de ruta por fases:

\subsection*{Fase 1 (0--3 meses): estabilización técnica}
\begin{itemize}
  \item Refactor de gestión de estado a Provider/BLoC.
  \item Cobertura mínima de tests unitarios para autenticación, rutinas y mensajería.
  \item Instrumentación básica de errores y rendimiento.
\end{itemize}

\subsection*{Fase 2 (3--6 meses): valor funcional diferencial}
\begin{itemize}
  \item Módulo de nutrición y seguimiento de adherencia.
  \item Notificaciones push por recordatorios de sesión y mensajes nuevos.
  \item Mejoras de UX orientadas a retención (onboarding, ayudas contextuales, estados vacíos guiados).
\end{itemize}

\subsection*{Fase 3 (6--12 meses): escalado y negocio}
\begin{itemize}
  \item Versión web de panel de entrenador.
  \item Integración de pagos y suscripciones.
  \item Despliegue controlado en tiendas con analítica de conversión.
\end{itemize}

\section{Impacto esperado de las mejoras}

La implantación de este roadmap tendría impacto en tres dimensiones:

\begin{enumerate}
  \item \textbf{Impacto técnico}: mayor robustez, mantenibilidad y reducción de incidencias en producción.
  \item \textbf{Impacto de producto}: incremento de utilidad real para entrenadores con necesidades de seguimiento completo.
  \item \textbf{Impacto económico}: transición de proyecto académico a producto monetizable mediante modelo de suscripción.
\end{enumerate}

Con este enfoque, FitLabs podría evolucionar de MVP académico a plataforma especializada con capacidad de adopción en entornos reales de entrenamiento personal.
